\begin{tabular}{@{}lrrrrrrp{3.9cm}@{}}
\hline
Materiale & $n$ & $k$ & $\rho$ & $\sigma_y$ & $T_m$ & $\Delta n$ & Note \\
 & & & (kg/m$^3$) & (GPa) & (K) & & \\
\hline
Cu lucido & 1{,}09 & 2{,}53 & 8960 & 0{,}2 & 1358 & 0 & assorbe il 40\% (Fresnel); $\varepsilon_{IR}$ da Mie 0,0045--0,0109: specchio nell'IR \\
Cu nero & 1{,}09 & 2{,}53 & 8960 & 0{,}2 & 1358 & 0 & $\alpha_{532}=0{,}95$; $\varepsilon_{IR}=0{,}30$ dichiarato (0,20--0,60), non Mie \\
Silice fusa & 1{,}46 & $1{,}00\times10^{-7}$ & 2200 & 1 & 1983 & 0 & trasparente; $\Delta n$ solo di forma: 0,02 disco, 0,03 manubrio, 0,01 anello/reticolo \\
Vaterite & 1{,}65 & $1{,}00\times10^{-6}$ & 2710 & 0{,}1 & 723 & 0{,}06 & $\Delta n=0{,}06$ efficace (0,06--0,10); $k=10^{-6}$; $T_m=723$ K (vaterite$\to$calcite); fragile, $\sigma_y$ conservativo \\
\hline
\end{tabular}

\par\smallskip\noindent\footnotesize La calcite ha $\Delta n=0{,}17$: e' il valore usato nel test E (Friese 1998), non la vaterite policristallina della mappa.
